% Image credits for the photographs and AI illustrations of Book 4 (University, Year 2).
% Machine-readable license records live in images/book4/CREDITS.md; keep the
% two files in sync. The Book 4 writer owns this file: add one \item per
% photograph (inside the itemize, by chapter) as photographs land.
\chapter*{Créditos de las imágenes}

Los dibujos, mapas y gráficos de este libro son originales; los mapas y gráficos
elaborados a partir de datos publicados citan su fuente en el pie de figura o en
esta lista. Las fotografías se usan bajo sus respectivas licencias libres, con
gratitud a sus autores; los enlaces completos a la fuente y las URL de las
licencias de cada fotografía figuran en \texttt{images/book4/CREDITS.md} en el
repositorio del libro.

\bigskip
Junto a las fotografías y a los dibujos originales, algunas figuras son
ilustraciones generadas por IA, producidas con la generación de imágenes de OpenAI
a partir de instrucciones escritas por los editores del libro, y revisadas por su
exactitud científica antes de incluirlas. Muestran solo escenas genéricas, nunca un
lugar real, una persona real ni un ejemplar real. La instrucción completa de cada
imagen generada figura en \texttt{images/book4/ai/PROMPTS.md} en el repositorio.

\bigskip
\noindent Fotografías y datos, por capítulo:
\begin{itemize}\small
  \item Cap.~1 (el medio interestelar y la formación de estrellas): el disco protoplanetario de HL Tauri, ALMA (ESO/NAOJ/NRAO), CC BY 4.0; los Pilares de la Creación en la nebulosa del Águila, NASA, ESA and the Hubble Heritage Team (STScI/AURA), dominio público.
  \item Cap.~2 (la formación de los planetas y los meteoritos): la condrita carbonácea de Allende y las figuras de Widmanstätten del meteorito metálico de Gibeon, James St.~John, CC BY 2.0.
  \item Cap.~3 (los elementos en la Tierra): xenolito del manto de lherzolita en un basalto, Kilbourne Hole, James St.~John, CC BY 2.0.
  \item Cap.~5 (los magmas y su emplazamiento): dique cerca de Shiprock, Dexter Perkins, CC0 (Wikimedia Commons).
  \item Cap.~6 (las erupciones volcánicas): fuente de lava en el cono Pu`u~`\=O`\=o del Kīlauea, J.\,D. Griggs, USGS, dominio público; la columna eruptiva del monte Santa Helena, 18 de mayo de 1980, Austin Post, USGS, dominio público; la columna eruptiva del monte Pinatubo, 12 de junio de 1991, Dave Harlow, USGS, dominio público.
  \item Cap.~7 (la fusión del manto): muestra de mano de dunita, Mike Davis (GeoDIL), CC0.
  \item Cap.~8 (la diferenciación magmática): capas de cromitita del complejo de Bushveld, Kevin Walsh, CC BY 2.0; gabros estratificados de la intrusión de Skaergaard, C. Tegner, CC BY 4.0.
  \item Cap.~10 (reacciones y facies metamórficas): esquisto azul de lawsonita, Ward Creek, California, y eclogita, Almenningen, Noruega, James St.~John, CC BY 2.0.
  \item Cap.~11 (termobarometría y trayectorias P--T--t): esquisto granatífero, James St.~John, CC BY 2.0; grano de coesita en una eclogita, con nicoles cruzados, J.\,R. Smyth, CC BY-SA 1.0.
  \item Cap.~12 (la fusión de la corteza: migmatitas y granitos): afloramiento de migmatita en Nashua, New Hampshire, Kevin M. Gill, CC BY 4.0; filón de leucogranito, Anne Burgess (Geograph), CC BY-SA 2.0.
  \item Cap.~13 (anatomía de una cordillera: los Alpes): lavas almohadilladas de la ofiolita del Chenaillet, Lionel Duchoiselle, CC BY-SA 3.0; el Cervino desde Zermatt, Ximonic (Simo Räsänen), CC BY-SA 4.0.
  \item Cap.~14 (cordilleras de colisión y cordilleras andinas comparadas): la meseta del Tíbet entre Nepal y Shigatse, Jan Kranendonk, CC BY-SA 4.0; el Altiplano boliviano, Luca Galuzzi, CC BY-SA 2.5.
  \item Cap.~15 (macizos antiguos y ciclo de Wilson): Siccar Point, acantilado sudeste, dave souza, CC BY-SA 4.0.
  \item Cap.~16 (cuencas sedimentarias comparadas): plataforma de perforación semisumergible en el mar del Norte, Erik Christensen, CC BY-SA 3.0.
  \item Cap.~17 (islas oceánicas): el atolón Kure desde la Estación Espacial Internacional, NASA, dominio público.
  \item Cap.~18 (la Tierra precámbrica): komatiita espinifex en Pyke Hill, Ontario, James St.~John, CC BY 2.0; carbonato de cobertera de la formación Nuccaleena, Australia Meridional, James St.~John, CC BY 2.0.
  \item Cap.~19 (la Tierra fanerozoica): pilares de creta de Monument Rocks, Kansas, James St.~John, CC BY 2.0; el nivel Bonarelli, Picenus, CC BY-SA 4.0.
  \item Cap.~20 (fósiles y paleoambientes): \emph{Waptia} y \emph{Marrella} de la lutita de Burgess, James St.~John, CC BY 2.0; el ejemplar de Berlín de \emph{Archaeopteryx}, H.~Raab, CC BY-SA 3.0; \emph{Zoophycos} en Sagres, Amada44, CC BY 4.0.
  \item Cap.~21 (termodinámica atmosférica): yunque de una tormenta supercelular, Stephen Corfidi, NOAA, dominio público; sondeos de radiosonda del Integrated Global Radiosonde Archive, I. Durre y otros, NOAA NCEI, dominio público.
  \item Cap.~22 (sistemas meteorológicos): un sistema de bajas presiones sobre Islandia, NASA/GSFC, MODIS Rapid Response Team, Jacques Descloitres, dominio público; el ojo del huracán Isabel desde la Estación Espacial Internacional, Ed Lu, NASA, dominio público.
  \item Cap.~23 (variabilidad climática): índice de la Oscilación del Atlántico Norte e Índice Oceánico de El Niño (ONI), NOAA Climate Prediction Center, dominio público.
  \item Cap.~24 (archivos e indicadores climáticos): capas anuales del testigo de hielo GISP2 a 1855~m, Anthony Gow, U.S. Army Corps of Engineers CRREL (NOAA Paleoclimatology), dominio público; sección de un espeleotema de travertino, Diamond Caverns, Kentucky, James St.~John, CC BY 2.0; el registro bentónico apilado LR04 de L.\,E. Lisiecki y M.\,E. Raymo (2005), el registro de temperatura de EPICA Dome C de J. Jouzel y otros (2007) y la serie compuesta antártica de \ce{CO2} de B. Bereiter y otros (2015), NOAA NCEI Paleoclimatology, dominio público.
  \item Cap.~25 (el forzamiento orbital y las glaciaciones): elementos orbitales e insolación de Berger y Loutre (1991), NOAA NCEI Paleoclimatology, dominio público.
  \item Cap.~26 (el clima a lo largo del tiempo geológico): registro isotópico bentónico del Cenozoico CENOGRID, T. Westerhold y otros (2020), PANGAEA, CC BY 4.0; compilación de \ce{CO2} del Fanerozoico de G.\,L. Foster, D.\,L. Royer y D.\,J. Lunt (2017), \emph{Nature Communications}, CC BY 4.0.
  \item Cap.~27 (el cambio climático antropógeno): Sexto Informe de Evaluación del IPCC, Grupo de Trabajo~I, datos ampliados del anexo~III, CC BY 4.0.
  \item Cap.~28 (peligros geológicos y riesgo): depósitos de flujos de derrubios en Caraballeda, Venezuela, Lawson Smith, U.S. Army Corps of Engineers, dominio público; una mezquita en pie tras el tsunami de 2004 en Aceh, USGS, dominio público.
  \item Cap.~29 (yacimientos minerales): la mina de Bingham Canyon desde la Estación Espacial Internacional, NASA (tripulación de la Expedición~15), dominio público; gossan de la mina Flambeau, Wisconsin, James St.~John, CC BY 2.0; datos de ley y tonelaje de D.\,A. Singer, V.\,I. Berger y B.\,C. Moring (2008), USGS, dominio público.
  \item Cap.~30 (geología de la energía): central geotérmica cerca de Geyserville, California, U.S. Department of Defense, dominio público; densidades del dióxido de carbono del NIST Chemistry WebBook (Standard Reference Database 69).
  \item Cap.~31 (métodos de campo y de laboratorio II): augita poiquilítica que engloba plagioclasa, Strekeisen, CC BY-SA 4.0; lámina delgada de un esquisto de granate y mica, Jackdann88, CC BY 3.0; microsonda electrónica del Servicio Geológico de Finlandia, Kallerna, CC BY-SA 4.0; resultados de sílice de materiales de referencia de rocas, J.\,E. Taggart y D.\,F. Siems (2002), USGS, dominio público.
\end{itemize}
